\documentclass[../../../include/open-logic-section]{subfiles}

\begin{document}

\olfileid{his}{set}{hilbertcurve}

\olsection{附录：Hilbert 的空间填充曲线}

在 \olref[pathology]{sec} 节中，我们提到 康托尔关于一条直线和一个正方形具有完全相同数量的点的证明（\olref[his][set][cantorplane]{thm:cantorplane}），促使 皮亚诺思考是否可能存在一种空间填充\emph{曲线}。他在 1890 年得到了肯定的答案。在本节中，我们将以略带示意性的方式解释如何构造 希尔伯特的空间填充曲线（并做一点微小的调整）。\footnote{更严格的解释，参见 \cite{Rose2010}。这一调整相当于在下文曲线中加入了红色部分。这使得验证该曲线的连续性稍容易一些。}

我们必须定义函数 $h$ 为函数序列 $h_1$, $h_2$, $h_3$, \dots 的极限。我们首先描述构造过程，然后证明它是空间填充的，最后证明它是一条曲线。

我们将取 $h$ 的值域为单位正方形 $\unitsquare$。以下是 $h$ 的第一个逼近，即 $h_1$：
\begin{center}
	\begin{tikzpicture}
	\draw[gray] (0pt, 0pt) rectangle (80pt, 80pt);
	\draw[step = 40pt, gray, very thin] (0pt, 0pt) grid (80pt, 80pt);
	\draw[oldiagcolorC, thick] (20pt, 0)--(20pt, 20pt);
	\draw[oldiagcolorC, thick] (60pt, 0)--(60pt, 20pt);		
	\begin{scope}[xshift=20pt, yshift=20pt]
	\draw [black, thick, l-system={Hilbert curve, step=40pt, angle=90, axiom=L, order=1}]  lindenmayer system;
	\end{scope}
	\end{tikzpicture}
\end{center}
为了看清构造过程，我们在正方形上叠加了一个 $2 \times 2$ 的网格。我们可以认为曲线起始于左下方四分之一区域，移动到左上方，再到右上方，最后到达右下方。以下是构造的第二阶段，即 $h_2$：
\begin{center}
	\begin{tikzpicture}
	\draw[gray] (0pt, 0pt) rectangle (80pt, 80pt);
	\draw[step = 20pt, gray, very thin] (0pt, 0pt) grid (80pt, 80pt);
	\draw[oldiagcolorC, thick] (0pt, 10pt)--(10pt, 10pt);
	\draw[oldiagcolorC, thick] (70pt, 10pt)--(80pt, 10pt);
	\draw[thick, oldiagcolorE] (10pt, 30pt)--(10pt, 50pt); 
	\draw[thick, oldiagcolorE] (30pt, 50pt)--(50pt, 50pt); 
	\draw[thick, oldiagcolorE] (70pt, 50pt)--(70pt, 30pt); \begin{scope}[xshift=10pt, yshift=30pt]
	\draw [thick, rotate=270,black, l-system={Hilbert curve, step=20pt, angle=90, axiom=L, order=1}]  lindenmayer system;
	\end{scope}
	\begin{scope}[xshift=10pt, yshift=50pt]
	\draw [thick, black, l-system={Hilbert curve, step=20pt, angle=90, axiom=L, order=1}]  lindenmayer system;
	\end{scope}
	\begin{scope}[xshift=50pt, yshift=50pt]
	\draw [thick, black, l-system={Hilbert curve, step=20pt, angle=90, axiom=L, order=1}]  lindenmayer system;
	\end{scope}
	\begin{scope}[xshift=70pt, yshift=10pt]
	\draw [thick, rotate=90,black, l-system={Hilbert curve, step=20pt, angle=90, axiom=L, order=1}]  lindenmayer system;
	\end{scope}
	\end{tikzpicture}
\end{center}
不同的颜色将有助于解释 $h_2$ 是如何构造的。我们首先将 $h_1$ 的非红色部分按比例缩小，分别放到正方形的左下、左上、右上和右下区域（用黑色绘制）。然后我们用绿色线条连接这四个图形。最后，我们用红色线条将图形连接到正方形的边界。

现在来看 $h_3$。正如 $h_2$ 是由四个按比例缩小的 $h_1$ 非红色部分的副本相连而成一样，$h_3$ 也是由四个按比例缩小的 $h_2$ 非红色部分的副本（用黑色绘制）相连而成，这些副本随后用绿色线条连接在一起，最后用红色线条连接到正方形的边界。
\begin{center}
	\begin{tikzpicture}
	\draw[gray] (0pt, 0pt) rectangle (80pt, 80pt);
	\draw[step = 10pt, gray, very thin] (0pt, 0pt) grid (80pt, 80pt);
	\draw[thick, oldiagcolorC](5pt, 0pt)--(5pt, 5pt); 
	\draw[thick, oldiagcolorC] (75pt, 0pt)--(75pt, 5pt); 
	\draw[thick, oldiagcolorE] (75pt, 45pt)--(75pt, 35pt);
	\draw[thick, oldiagcolorE] (5pt, 35pt)--(5pt, 45pt); 
	\draw[thick, oldiagcolorE] (35pt, 45pt)--(45pt, 45pt); 
	\draw[thick, oldiagcolorE] (75pt, 45pt)--(75pt, 35pt); \begin{scope}[xshift=5pt, yshift=35pt]
	\draw [rotate=270, thick, black, l-system={Hilbert curve, step=10pt, angle=90, axiom=L, order=2}]  lindenmayer system;
	\end{scope}
	\begin{scope}[xshift=5pt, yshift=45pt]
	\draw [thick, black, l-system={Hilbert curve, step=10pt, angle=90, axiom=L, order=2}]  lindenmayer system;
	\end{scope}
	\begin{scope}[xshift=45pt, yshift=45pt]
	\draw [thick, black, l-system={Hilbert curve, step=10pt, angle=90, axiom=L, order=2}]  lindenmayer system;
	\end{scope}
	\begin{scope}[xshift=75pt, yshift=5pt]
	\draw [thick, rotate=90,black, l-system={Hilbert curve, step=10pt, angle=90, axiom=L, order=2}]  lindenmayer system;
	\end{scope}
	\end{tikzpicture}
\end{center}
现在，我们看到了从 $h_n$ 定义 $h_{n+1}$ 的一般模式。
最后我们通过考虑这些相继函数 $h_1$, $h_2$,
\dots\ 的逐点极限来定义曲线 $h$\emph{本身}。
也就是说，对每个 $x \in \unitsquare$：
\begin{align*}
	h(x) &= \lim_{n \rightarrow \infty} h_n(x)
\end{align*} 
现在我们证明这条曲线可以填满空间。当我们画曲线 $h_n$ 时，
我们在 $\unitsquare$ 上叠加了一个 $2^n \times 2^n$ 的网格。
根据勾股定理，每个网格位置的对角线长度为：
\[
\sqrt{\left(\nicefrac{1}{2^{n}}\right)^2+\left(\nicefrac{1}{2^{n}}\right)^2} = 2^{(\frac{1}{2}-n)}
\]
并且显然 $h_n$ 经过了每一个网格位置。所以 $\unitsquare$ 中的每个点
与 $h_n$ 上的某点距离\emph{至多}为 $2^{(\frac{1}{2}-n)}$。
现在，$h$ 被定义为函数 $h_1$, $h_2$, $h_3$, \dots\ 的极限。
因此，任一点到 $h$ 的最大距离由下式给出：
\[
\lim_{n \rightarrow \infty} 2^{(\frac{1}{2}-n)} = 0.
\]
这就是说：$\unitsquare$ 中的每一点到~$h$ 的距离为 $0$。
换言之，$\unitsquare$ 的每一点都位于这条曲线上。所以 $h$ 填满了空间！{}

剩下的问题是证明 $h$ 确实是一条\emph{曲线}。为此，
我们必须先定义这个概念。现代定义基于 Jordan（若尔当）在1887年给出的一个定义
（即在第一条空间填充曲线被提出前仅几年）：

\begin{defn}
曲线是从 $\unitline$ 到 $\Real^2$ 的连续映射。
\end{defn}

这一定义相当直观：粗略地说，曲线就是将一条典范线映射到平面 $\Real^2$ 上的“平滑”映射。
我们的函数 $h$ 确实是一个从 $\unitline$ 到 $\Real^2$ 的映射。
所以，我们只需证明 $h$ 是连续的。
我们在 \olref[limits]{sec} 中用 $\epsilon$/$\delta$ 记号定义了连续性。
通俗地说，我们要确立以下事实：\emph{如果你在 $\unitsquare$ 中指定一点 $p$，
以及任意想要的精度 $\epsilon$，我们都能找到 $\unitline$ 的一个开区间，
使得对于该开区间中的任意 $x$，$h(x)$ 都在 $p$ 的 $\epsilon$ 范围内。}

那么：假设你已经指定了 $p$ 和 $\epsilon$。
这实际上相当于在 $\unitsquare$ 上画一个以 $p$ 为中心、$\epsilon$ 为半径的圆。
（这个圆可能会超出 $\unitsquare$ 的边界，但这没关系。）
现在回想一下，在描述函数 $h_n$ 时，我们在 $\unitsquare$ 上叠加了一个 $2^n \times 2^n$ 的网格。
显然，不论 $\epsilon$ 多小，总存在某个 $n$，使得 $\unitsquare$ 上 $2^n \times 2^n$ 网格的
某个单独的网格位置完全落在以 $p$ 为中心、$\epsilon$ 为半径的圆内。

所以，取这个 $n$，并令 $I$ 为 $\unitline$ 上使得 $h_n$ 完全映射到相关网格位置的最大开区间。
（显然这样的 $(a,b)$ 是存在的，因为我们已经注意到 $h_n$ 经过了 $2^n\times 2^n$ 网格中的每一个网格位置。）
现在只需证明，只要 $x \in I$，点 $h(x)$ 就位于同一个网格位置中。
而要证明\emph{这一点}，只需证明对任意 $m > n$，$h_m(x)$ 都位于同一个网格位置中。
但这是显而易见的。如果我们考虑 $m > n$ 时 $h_m$ 的情况，
会发现单位区间的“相同部分”被映射到了同一个网格位置；
我们只是以一种越来越拉伸、蜿蜒的方式将其映射到该区域。

\end{document}
